La distribución de temperatura en una placa metálica cuadrada $[-2,2]\times[-2,2]$ está modelada por:
\[ T(x,y) = 100 - 10x^2 - 20y^2 \quad (^\circ\text{C}) \]
\begin{enumerate}
    \item Encontrar el dominio y la imagen de $T$ sobre la placa.
    \item Calcular la temperatura en los puntos $(0,0)$, $(1,1)$ y $(2, 0)$.
    \item Trazar las curvas de nivel (isotermas) para $T = 100, 80, 60, 40, 20$ °C. ¿Qué forma tienen? ¿Cuáles de ellas intersectan la placa?
    \item Determinar el conjunto de puntos de la placa donde la temperatura es de exactamente $70$ °C e identificarlo geométricamente.
    \item Una partícula se mueve sobre la placa siguiendo la trayectoria $\vec{r}(t) = (t, t^2 - 1)$ para $t \in [-1, 1]$. Calcular $T(x(t), y(t))$ a lo largo de esta trayectoria y encontrar el punto de temperatura máxima.
\end{enumerate}